\section{Folgen und Reihen}

\Title{Definitionen}

\Def Eine \textbf{Folge} (reeller Zahlen) ist eine Abbildung $a: \N^{*} \rightarrow \R $\\
Wir schreiben $a_{n}$ statt $a(n)$ und bezeichnen eine Folge mit $(a_{n})_{n \geq 1}$\\
\Bsp $a_{n} = \frac{1}{n}$

\Def Sei $(a_{n})_{n \geq 1}$ eine Folge in $\R $ oder $\C$, so ist die \textbf{Reihe} 
$(S_{n})_{n \geq 1}$, die Partialsumme $S_n = \sum_{k=1}^{n}a_k$.


\Title{Wichtige Folgen}

\begin{enumerate}[label=(\arabic*)]
	\item{Arithmetische Folge: \((a_n)\), sodass \(a_{n+1}-a_n=d\) konstant ist, d.h. \(a_1=a,\;a_2=a_1+d,\quad a_n=a+(n-1)d\)}
	\item{Geometrische Folge: \((a_n)\), sodass \(a_{n+1}=qa_n\) mit \(q\) konstant, d.h. \(a_1=a,\;a_2=qa,\quad a_n=q^{n-1}\cdot a\)}
\end{enumerate}

\Satz Es gilt \(a+qa+\dots+q^{n-1}a = a(1+q+\dots+q^{n-1}) = \begin{cases} na&\text{für }q=1\\ a\cdot\frac{1-q^n}{1-q}& \text{für }q\neq 1\end{cases}\)


\Title{Grenzwert und Konvergenz}

\begin{tcolorbox}
	\Def Eine Folge $(a_n)_{n \geq 1}$ in $\C$ \textbf{konvergiert} für $n \rightarrow \infty$ gegen ein
	$a \in \C$, falls es für alle $\epsilon > 0$ ein $N \in \N$ gibt, sodass für alle 
	$n \geq \N$ gilt $|a_n-a| < \epsilon$. \\

	\textbf{Äquivalent:} Eine Folge $(a_n)_{n \geq 1}$ heisst konvergent, falls es $l \in \R$ gibt,
	so dass $\forall \epsilon > 0$ die Menge $\{n \in \N^{*}: a_{n} \notin ]l - \epsilon, l + \epsilon[\}$ endlich ist.
	Man bezeichnet $a= \lim_{n \rightarrow \infty} a_n$, wobei $a$ der Grenzwert von $(a_n)$ ist. \\

	Falls die Reihe $(S_n)_{n \in \N}$ gegen $S \in \C$ konvergiert, bezeichnet man $S=\sum_{n=1}^\infty a_n$
\end{tcolorbox}

Falls eine Folge oder Reihe nicht konvergiert, so divergiert sie. Schreiben wir $a=\lim_{n\rightarrow\infty}a_n$, 
so meinen wir, dass der Grenzwert der Folge $(a_n)$ existiert und gleich $a$ ist.

\Bsp (1) Sei $k \in \N$, $a_n = n^kq^n$, $0\leq|q|<1 \implies \lim_{n\rightarrow\infty}a_n = 0\)\\
\hspace{8mm}(2) Sei $c \geq 1$, $a_n=\frac{c^n}{n!} \implies \lim_{n\rightarrow\infty}a_n = 0\) \\
\hspace{8mm}(3) Sei b $\in \mathbb{Z},$ $\lim_{n \to \infty} (1 + \frac{1}{n})^{b}$ = 1


\begin{tcolorbox}
	\Satz Eine \textbf{konvergente Folge ist immer beschränkt} und hat \textbf{genau} einen Grenzwert.
\end{tcolorbox}

\color{blue}
\Bem konvergent $\Rightarrow$ beschränkt, aber nicht umgekehrt!	\\
\Bem $a_n \quad konv. \Rightarrow b_n = a_{n+1}-a_n \quad konv$
\color{black}

Um zu zeigen, dass $a$ der Grenzwert der Folge $a_n$ ist, müssen wir zeigen, dass $|a_n - a| < \epsilon$
ab einem gewissen $n$ für alle $\epsilon$ gilt. D.h. wir formen nach $n$ um, dabei müssen wir auf das Ungleichzeichen
und mögliche Wurzeln aufpassen.

\begin{tcolorbox}
	\Satz (Sandwich-Theorem)
	Sei $\Limit a_n = \alpha$ und $\Limit c_n = \alpha$ und $a_n \leq b_n \leq c_n, \forall n \geq k$,
	dann gilt $\Limit b_n = \alpha$
\end{tcolorbox}


\vspace*{0mm}
\columnbreak
%----------------------------------------------------------------------------------------------------


\Satz Falls $(b_n)_{n \geq 1}$ gegen 0 konvergiert und \(|a_n-a|\leq b_n\), dann konvergiert 
\((a_n)_{n \geq 1}\) gegen \(a\).

\Satz Seien $(a_n)_{n \geq 1},(b_n)_{n \geq 1}$ mit $a = \lim_{n \rightarrow \infty} a_n$ und 
$b = \lim_{n \rightarrow \infty} b_n$, so gilt:
\begin{enumerate}[label=(\arabic*)]
    	\item{Die Folge \((a_n+b_n)_{n \geq 1}\) konvergiert gegen \(a+b\)}
    	\item{Die Folge \((a_n\cdot b_n)_{n \geq 1}\) konvergiert gegen \(a\cdot b\)}
    	\item{Falls $b \neq 0$ ist $b_n \neq 0$ für $n$ gross genug und 
    	$\lim_{n \rightarrow \infty} \left( \frac{a_n}{b_n} \right) = \frac{a}{b}$.}
    	\item{Falls $\exists K \geq 1$ sodass \(a_n\leq b_n\), $\forall n \geq K$, ist \(a\leq b\)}
\end{enumerate}


\Title{Der Satz von Weierstrass}

\Def $(a_{n})_{n \geq 1}$ ist monoton wachsend [fallend] falls: $a_{n} \leq a_{n+1} 
[a_{n} \geq a_{n+1}], \forall n \geq 1$.

\begin{tcolorbox}
	\Satz Sei ($a_{n})_{n \geq 1}$ monoton wachsend und nach oben beschränkt. Dann konvergiert $(a_{n})_{n \geq 1}$ mit Grenzwert
			$\lim_{n \to \infty} a_{n} = \sup\{a_{n} : n \geq 1\} $ \\
		
		Sei ($a_{n})_{n \geq 1}$ monoton fallend und nach unten beschränkt. Dann konvergiert $(a_{n})_{n \geq 1}$ mit Grenzwert
			$\lim_{n \to \infty} a_{n} = \inf\{a_{n} : n \geq 1\} $
\end{tcolorbox}


\Title{Bolzano Weierstrass}

\Def Eine \textbf{Teilfolge} einer Folge $(a_{n})_{n \geq 1}$ 
ist eine Folge $(b_{n})_{n \geq 1}$ mit $b_{n} = a_{l(n)}$, wobei $l: 
N^{*} \rightarrow N^{*}$ eine Abbildung bezeichnet mit der Eigenschaft $l(n) < l(n+1)$ 
$\forall n \geq 1$ (das jeweils nächste Element der Teilfolge kommt auch später in der eigentlichen Folge) \\ \smallskip
\Satz Jede beschränkte Folge besitzt eine konvergente Teilfolge.


\Title{Limes Superior und Limes Inferior}

Mit jeder beschränkten Folge $(a_{n})_{n \geq 1}$ lassen sich zwei monotone Folgen $(b_{n})_{n \geq 1}$ 
und $(c_{n})_{n \geq 1}$ definieren, welche dann einen Grenzwert besitzen. \\
Sei für jedes n $\geq 1$: 
$$b_{n} = \inf\{a_{k} : k \geq n\} \text{\quad und \quad} c_{n} = \sup\{a_{k}: k \geq n\}$$

Dann ist $b_{n} \leq b_{n+1}$ und $c_{n+1} \leq c_{n}$ $\forall n \geq 1$. Nach 
Weierstrass sind beide folgen Konvergent, und wir definieren:
$$\liminf_{n \to \infty} a_{n} := \lim_{n \to \infty} b_{n} \; \text{ und } \; \limsup_{n \to \infty} a_{n} := \lim_{n \to \infty} c_{n}$$
Ist $\liminf_{n \to \infty} a_{n} = \limsup_{n \to \infty} a_{n}$ $\Rightarrow \exists \lim_{n \to \infty} a_{n}$ 
mit $\lim_{n \to \infty} a_{n} = \liminf_{n \to \infty} a_{n} = \limsup_{n \to \infty} a_{n}$

\Def Ein \textbf{abgeschlossenes Intervall} ist von der Form $[a,b]$, $]-\infty, b]$, $[a,\infty[$
oder $]-\infty, \infty[$. Sei $I = [a,b]$, so ist die \textbf{Länge} des Intervalls L$(I)$ = $b - a$. 

\Satz (Cauchy-Cantor) \\
Sei $I_{1} \supseteq I_{2} \supseteq I_{3} \supseteq ... I_{n} \supseteq I_{n+1} \supseteq...$ 
eine Folge abgeschlossener Intervalle mit $L(I_{1}) < + \infty$, dann gilt 
$\cap_{n \geq 1} I_{n} \neq \emptyset$. Ist $\lim_{n \to \infty}$ $L_{I_{n}} = 0$, so enthält 
die Schnittmenge genau einen Punkt.


\Title{Konvergenzkriterien}

\Bem Sei $(a_n)_{n \geq 1}$ konvergent $\Leftrightarrow$ $(a_n)_{n \geq 1}$ beschränkt
und $\liminf_{n \rightarrow \infty} a_n = \limsup_{n \rightarrow \infty} a_n$.

\Satz Falls $(a_n)_{n \geq 1}$ eine monoton wachsende [fallende] Folge ist, dann konvergiert $(a_n)_{n \geq 1}$ 
und $\lim_{n\rightarrow\infty}a_n=\Sup_{n \geq 1}(a_n)$ $[\Inf_{n \geq 1}(a_n)]$ genau dann, wenn $(a_n)_{n \geq 1}$
beschränkt ist.

Eine Folge in $\R^{d}$ verhält sich eigentlich gleich wie Folgen in $\R$. Sei 
$a(_{n})_{n \geq 1}$(gleiche Schreibweise) eine Folge in $R^{d}$.
Sie konvergiert falls: \\ 
$\exists a \in \R^{d}$ mit $\forall \epsilon > 0$ $\exists N \geq 1$ mit $||a_{n} - a|| < \epsilon$ $\forall n \geq N$


\vspace*{0mm}
\columnbreak
%----------------------------------------------------------------------------------------------------


\Title{Das Cauchy-Kriterium}

\begin{tcolorbox}
	\Satz Eine Folge $(a_n)_{n \geq 1}$ ist genau dann konvergent, falls
	$$\forall \epsilon > 0, \exists N \geq 1 \text{\quad so dass \quad} |a_n - a_m| < \epsilon \quad \forall n,m \geq N.$$
	Solch eine Folge nennt man auch \textbf{Cauchy-Folge}.
\end{tcolorbox}

\Bem Für eine Cauchy-Folge $(a_n)_{n \geq 1}$ gilt:
\begin{enumerate}[label=(\arabic*)]
	\item $\Rightarrow (a_n)_{n \geq 1}$ ist beschränkt
	\item $\Leftrightarrow (a_n)_{n \geq 1}$ ist konvergent 
\end{enumerate}

\Bem $a_n = \sum_{i=1} \frac{1}{i}$ konvergiert nicht, aber \(|a_{n+1}-a_n|\) konvergiert. 


\Title{Strategie - Konvergenz von Folgen}
\begin{enumerate}[label=(\arabic*)]
	\item Für \textbf{Brüche}, grösste Potenz von $n$ ausklammern und kürzen. Alle übrigen Brüche der Form $\frac{a}{n^s}$ streichen,
	da diese zu 0 konvergieren.
	\item Für \textbf{Wurzeln in einer Summe}, multipliziere mit der Differenz der Summe (bei $a + b$
	multipliziere mit $a - b$). 	
	\item Anwendung \textbf{Sandwich-Theorem}.
	\item Vergleich mit Referenz-Folgen.
	\item Grenzwert durch simple Operationen und Umformen ermitteln.
	\item Definition der Konvergenz / Limes anwenden.
	\item Für \textbf{rekursive Folgen}, Satz von Weierstrass anwenden, obere / untere Schranke abschätzen
	und per Induktion beweisen.
	\item Anwendung des Cauchy-Kriteriums
	\item Suchen eines konvergenten Majoranten.
\end{enumerate}

\Bsp $d_1 = 3, \; d_{n+1} = \sqrt{3d_n -2}, \; \forall n > 1$ \\
Beweise per Induktion, dass die Folge monoton fallend ist. Nun brauchen wir eine untere Schranke.
Induktionstrick für mögliche Kandidaten vom Limes. Angenommen die Folge konvergiert, so hat jede
Teilfolge den gleichen Grenzwert. Betrachten wir die Teilfolge $l(n) = n+1$:
$$d = \lim\limits_{n \to \infty} d_n =\lim\limits_{n \to \infty} d_{n+1} = \sqrt{\lim\limits_{n \to \infty} 3d_n -2} = \sqrt{3d -2}$$
$d^2 = 3d -2$, nehmen wir $d = 2$. Zeige nun dass die Folge nach unten beschränkt ist, mit d = 2.
Dies erfolgt wieder per Induktion.

\Title{Strategie - Divergenz von Folgen}

\begin{enumerate}[label=(\arabic*)]
	\item Suche einen divergenten Minoranten
	\item Für alternierende Folgen zeige, dass $\Limit a_{p_1(n)} \neq \Limit a_{p_2(n)}$.
\end{enumerate}


\Title{Limes Binom Trick}

Gegeben die Summe von zweier (oder einer) Wurzel, kann man wie folgt vorgehen:
$$\Lim_{x \to \infty} (\sqrt{x + 5} - \sqrt{x - 3}) = \Lim_{x \to \infty} (\frac{(x + 5) - (x - 3)}{\sqrt{x + 5} + \sqrt{x - 3}})$$


\Title{Limes Substitution Trick}

Hier ein Beispiel:
$$\Lim_{x \to \infty} x^2(1 - \cos(\frac{1}{x}))$$
Substituiere nun $u = \frac{1}{x}$:
$$\Lim_{u \to 0} \frac{1 - \cos(u)}{u^2} = \Lim_{u \to 0} \frac{\sin(u)}{2u} = \Lim_{u \to 0} \frac{\cos(u)}{2} = \frac{1}{2}$$


\vspace*{0mm}
\columnbreak
%----------------------------------------------------------------------------------------------------


\Title{Limes Log Trick}

Limes der Form $\infty^0$ und $1^\infty$ können meist mit dem Log-Trick berechnet werden:
$f(x)^{g(x)} = e^{g(x) \cdot \ln(f(x))}$, dann Bernoulli ($e$ ist stetig, daher betrachten wir nur den
 Exponenten) anwenden oder vereinfachen.

$\lim u(x)^{v(x)} = \lim e^{v(x)\cdot ln(u(x))} = e^{\lim [v(x)\cdot ln(u(x))]} = \begin{cases}
	u^v\\
	+\infty\\
	-\infty
\end{cases}$


\Title{Rechnen mit Reihen}

\begin{tcolorbox}
	\Satz Seien $\sum_{k=1}^{\infty} a_{k}$ und $\sum_{j=1}^{\infty} b_{j}$ konvergent, sowie $\alpha \in C$
	\begin{enumerate}[label=(\arabic*)]
		\item Dann ist $\sum_{k=1}^{\infty} (a_{k} + b_{k})$ = $(\sum_{k=1}^{\infty} a_{k}) + (\sum_{j=1}^{\infty} b_{j})$ und konvergiert.
		\item Dann ist $\sum_{k=1}^{\infty} (\alpha \cdot a_{k}) = \alpha \cdot \sum_{k=1}^{\infty} a_{k}$ und konvergiert.
	\end{enumerate}
\end{tcolorbox}


\Title{Cauchy-Kriterium für Reihen}

\Satz Die Reihe $\Reihe a_k$ konvergiert genau dann, falls
\vspace*{-6pt}
$$\forall \epsilon > 0, \exists N \geq 1 \text{\quad mit \quad} |\sum_{k = n}^{m}a_k| < \epsilon \quad \forall m \geq n \geq N$$

\begin{tcolorbox}
	\textbf{(Nullfolgenkriterium)} Aus dem Cauchy-Kriterium folgt, dass eine Reihe $\Reihe a_n $ divergiert, 
	falls $\Limit a_n \neq 0$.
\end{tcolorbox}


\Kor Sei $a_n \geq 0,\; s_n = a_1 + \dots + a_n$. So gilt:
\vspace{-6pt}
\begin{center}
	$\sum_{n=1}^\infty a_n \text{ konvergiert } \Longleftrightarrow (s_n) \text{ ist v.o.b. }$
\end{center}
\vspace{-6pt}
Insbesondere: Falls \(0\leq b_n\leq a_n\), so konvergiert auch $\sum_{n=1}^\infty b_n$.
Falls \(a_n\geq0\) und \(\sum_{n=1}^\infty a_n\) divergent ist, sagen wir \(\sum_{n=1}^\infty a_n=\infty\).

\Satz Sei $\sum_{k=1}^{\infty} a_{k}$ eine Reihe mit $a_{k} \geq 0,$ $\forall k \in \N^{*}$. Die Reihe $\sum_{k=1}^{\infty} a_{k}$ konvergiert genau dann, falls die Folge $(S_{n})_{n\geq 1}$, $S_{n} = \sum_{k=1}^{n} a_{k}$
der Partialsummen nach oben beschränkt ist.


\Title{Absolute Konvergenz}

\begin{tcolorbox}
	\Def Eine Reihe $\Reihe a_k, \; a_n \in \C$ heisst \textbf{absolut konvergent}, falls
	$\Reihe |a_k|$ konvergiert. \smallskip \\
	\Satz Eine absolut konvergente Reihe ist auch konvergent und es gilt:
	\vspace{-2mm}
	$$\left| \Reihe a_k \right| \leq \Reihe |a_k|$$
\end{tcolorbox}


\Title{Konvergenzkriterien}

\begin{tcolorbox}
	\Kor (Vergleichssatz)\\
	Seien $\Reihe a_{k}$ und $\Reihe b_{k}$ Reihen mit $0 \leq a_{k} \leq b_{k}$ $\forall k \geq K \geq 1$ 
	Dann gilt:
	\begin{enumerate}[label=(\arabic*)]
		\item $\Reihe b_{k}$ konvergent $\Rightarrow$  $\Reihe a_{k}$ konvergent
		\item $\Reihe a_{k}$ divergent $\Rightarrow$  $\Reihe b_{k}$ divergent
	\end{enumerate}
	Dies wird auch \textbf{Majoranten- / Minorantenkriterium} genannt.
\end{tcolorbox}

\begin{tcolorbox}
	\Satz (Leibnitz) \\
	Sei $(a_{n})_{n\geq 1}$ monoton fallend mit $a_{n} \geq 0,$ $\forall n \geq 1$ und $\Limit a_{n} = 0$. 
	Dann konvergiert
	\vspace{-1mm}
	\begin{center}
		S:= $\sum_{k=1}^{\infty}(-1)^{k+1} a_{k}$ 
		und es gilt: \textbf{$a_{1} - a_{2} \leq S \leq a_{1}$}.
	\end{center}
	\vspace{-1mm}
	Restgliedabschätzung: $|\sum_{k= n+1}^{\infty} (-1)^{k+1} a_k| \leq a_{n+1}$
\end{tcolorbox}

\color{blue}
\Satz (Dirichlet) Falls eine Reihe absolut konvergiert, dann konvergiert auch jede Umordnung
der Reihe mit dem selben Grenzwert. Eine Umordnung ist eine bijektive Abbildung $\phi : \N^{*} \to \N^{*}$.

\Satz (Riemann) Fall die Reihe nur konvergiert, so gibt es immer eine Annordnung, so 
dass $\Reihe a_{\phi(k)} = x, \forall x \in \R$.
\color{black}

\begin{tcolorbox}
	\textbf{Quotientenkriterium}: Falls $(a_n)_{n \geq 1}, a_n \neq 0$ und $\limsup_{n \to \infty}$ $\left|a_{n+1} / {a_n}\right|<1$, 
	dann konvergiert die Reihe abs. falls $\liminf_{n \to \infty}$ $\left|a_{n+1} / {a_n}\right|>1$ so divergiert sie.
	Äquivalent: Falls $\Limit\left|a_{n+1} / {a_n}\right|<1$, so ist die Reihe absolut konvergent (für \(>1\) divergent). \\

	\textbf{Wurzelkriterium}: Falls $\limsup_{n \to \infty} \sqrt[n]{|a_n|}<1$ konvergiert die Reihe absolut,
	falls $>1$ (auch $\limsup$) so divergiert sie.
	Äquivalent: Falls $\lim_{n\rightarrow\infty}\sqrt[n]{|a_n|}<1$, so ist die Reihe absolut konvergent (und für $>1$ divergent).
\end{tcolorbox}


\Bsp Sei \(0\neq z\in\C\) und \(a_n=\frac{z^{n}}{(n)!},\;n\in\N\).\\
\(\qquad \left|\frac{a_{n+1}}{a_n}\right|=\frac{|z|^{(n+1)}}{(n+1)!} {\Bigg/} \frac{|z|^{n}}{(n)!} = \frac{|z|}n \leq \frac 12<1 \) für \(n>2|z|\).\\
\(\qquad\implies\sum_{n=0}^\infty \frac{z^{n}}{(n)!}\) ist konvergent.


\Title{Potenzreihen}

Eine Potenzreihe ist eine Reihe von der Form $\sum_{k = 0}^{\infty} c_k \cdot z^k$, wobei $(c_k)_{k \geq 0}$
eine Folge ist. 

\begin{tcolorbox}
	\Def Der Konvergenzradius $\rho$ einer Potenzreihe entspricht:\\
	$\rho = 
	\begin{cases}
		+\infty & \text{falls } \limsup_{k\rightarrow\infty}\sqrt[k]{|c_k|}=0 \\
		\frac{1}{\limsup_{k\rightarrow\infty}\sqrt[k]{|c_k|}} & \text{falls }\limsup_{k\rightarrow\infty}\sqrt[k]{|c_k|}>0
	\end{cases}$

	(Äquivalent mit dem Quotientenkriterium)
\end{tcolorbox}

\Kor Die Potenzreihe $\sum_{k=0}^{\infty} c_{k}z^{k}$ konvergiert absolut für alle $|z| < \rho$ und 
divergiert für alle $|z| > \rho$. Der Fall $|z| = \rho$ muss jeweils noch separat geprüft werden.

\Satz Ein Potenzreihe $f(x) = \Reihe c_k x^k$ mit positivem Konvergenzradius $\rho$, konvergiert gleichmässig 
auf $[-\rho , \rho]$, insbesondere ist $f : ]-\rho, \rho [ \to \R$ stetig.

\Bem Absolut konvergente Potenzreihen sind stetig.

\Title{Doppelreihen}

Gegeben eine Doppelfolge $(a_{ij})_{i,j \geq 0}$, so können $\sum_{i = 0}^{\infty} \sum_{j = 0}^{\infty} a_{ij} = S_0 + S_1 + ...$
und $\sum_{j = 0}^{\infty} \sum_{i = 0}^{\infty} a_{ij} = b_0 + b_1 + ...$ beide konvergent sein mit
verschiedenen Grenzwerten. Wir nennen $\sum_{i,j \geq 0} a_{ij}$ eine Doppelreihe. Wenn die Reihe absolut
konvergiert, so sind beide Grenzwerte gleich und jede Anordnung konvergiert zum selben Grenzwert.


\vspace*{0mm}
\columnbreak
%----------------------------------------------------------------------------------------------------


\Title{Das Cauchy Produkt}

\begin{tcolorbox}
	\Def Das Cauchy Produkt zweier Reihen $\sum_{i = 0}^{\infty} a_i$ und $\sum_{j = 0}^{\infty} b_j$ ist
	die Reihe: $$\sum_{n = 0}^{\infty} \sum_{j = 0}^{n} (a_{n-j} \cdot b_j) = a_0b_0 + (a_0b_1 + a_1b_0) + ...$$ 
	Falls beide Reihen konvergieren, so konvergiert auch das Cauchy Produkt.
\end{tcolorbox}
	

Nun wollen wir noch betrachten ob man Summation und Limes vertauschen kann.

\Satz Sei $f_{n}: \N \rightarrow \R$ eine Folge mit:
\begin{enumerate}
	\item {$f(j) := \lim_{n \to \infty} f_{n}(j)$ existiert $\forall j \in \N$}
	\item{Es gibt eine Funktion g: $\N \to [0,\infty[$, mit \\
		$|f_{n}(j)| \leq g(j), \; \forall j \geq 0, \forall n \geq 0$ und 
		$\sum_{j=0}^{\infty} g(j)$ konvergiert. 
	}
\end{enumerate}
Dann gilt $\sum_{j=0}^{\infty} f(j) = \lim_{n \to \infty} \sum_{j=0}^{\infty} f_{n}(j)$.


\Title{Integral Test}

\begin{tcolorbox}
	Sei $f(x)$ eine stetige, positive und monoton fallende Funktion auf $[k, \infty[$ und $f(n) = a_n$:
	$$\int_k^\infty f(x)dx \text{ konvergiert } \Leftrightarrow \Sum_{n = k}^{\infty} a_n \text{ konvergiert }$$
	$$\int_k^\infty f(x)dx \text{ divergiert } \Leftrightarrow \Sum_{n = k}^{\infty} a_n \text{ divergiert }$$
\end{tcolorbox}


\Title{Strategie - Konvergenz von Reihen}

\begin{enumerate}[label=(\arabic*)]
	\item Handelt es sich um eine spezielle Reihe? (Geometrisch, Teleskopiert, Harmonisch, Zetafunktion)
	\item Ist $\Limit a_n = 0$? (Nullfolgenkriterium)
	\item Ist das Quotientenkriterium oder Wurzelkriterium anwendbar?
	\item Existiert ein konvergierender Majorant / divergirender Minorant?
	\item Kann man das Leibnitzkriterium anwenden?
	\item Integral Test?
\end{enumerate}
